\renewcommand{\e}{\vec{e_1}}
\newcommand{\ee}{\vec{e_2}}

\renewcommand{\u}{\vec{u_1}}
\newcommand{\uu}{\vec{u_2}}

\renewcommand{\v}{v_1}
\newcommand{\vv}{v_2}

\newcommand{\vp}{v_1^\prime}
\newcommand{\vvp}{v_2^\prime}
\newcommand{\hpi}{\frac{-\pi}{2}}


\begin{frame}
	\begin{columns}[t]
		\begin{column}{0.6\textwidth}
			\textbf{Changement de base:}
			\quickfigure{\footnotesize{Changement de base}}{base-change
			}{.66}{base-change.png}\vspace{-.7cm}

			\footnotesize{\begin{align}
					E &= \{ \e, \ee \}
					&U = \{ \u, \uu \} \\
					\e &= 0\u + (-1)\uu
					&\ee = \u + 0\uu
			\end{align}
			\begin{align}
					\v\e + \vv\ee &= \vp\u + \vvp\uu \\
					\v(0\u - \uu) + \vv(\u + 0\uu) &= \vp\u + \vvp\uu \\
					(0\v + \vv)\u + (-\v + 0\vv)\uu &= \vp\u + \vvp\uu \\
					\begin{bmatrix} 0 & 1 \\ -1 & 0 \end{bmatrix}
					\begin{bmatrix} \v \\ \vv \end{bmatrix} &=
					\begin{bmatrix} \vp \\ \vvp \end{bmatrix}
			\end{align}}


		\end{column}
		\begin{column}{0.4\textwidth}
		\textbf{Matrice de passage:}

		\footnotesize{\begin{align}
		\theta &= \hpi \\
		T &= \begin{bmatrix}
						\cos\theta & -\sin\theta \\
						\sin\theta & \cos\theta \\
         \end{bmatrix} \\
		T &= \begin{bmatrix}
						\cos\hpi & -\sin\hpi \\
						\sin\hpi & \cos\hpi \\
         \end{bmatrix} \\
		T &= \begin{bmatrix}
						0 & 1 \\
						-1 & 0 \\
         \end{bmatrix}
		\end{align}}

		\quickfigure{
			\tiny{Exemple de rotation}
		}{rotated-example}{.6}{rotated-example.png}

		\end{column}
	\end{columns}
\end{frame}
